\documentclass[letterpaper, 12pt]{article}

\usepackage[spanish]{babel}
\usepackage[utf8]{inputenc}
\usepackage[T1]{fontenc}
\usepackage{graphicx}

\pagestyle{headings}

\newcommand{\portada}[1][documentTitle]{
    \begin{titlepage}
	    \centering
    	\includegraphics[width=0.30\textwidth]{src/img/escom}\par\vspace{1cm}
    	{\LARGE Instituto Politécnico Nacional \par}
	    {\Large Escuela Superior de Cómputo\par}
	    \vspace{1.5cm}
	    {\huge\bfseries #1 \par}
	    \vspace{2cm}
	    {\Large\itshape Valentin Ramos Emmanuel Guadalupe\par}
	    \vfill
	    Grupo: 1CM16\par
	    \vfill
	    {\large\today\par}
    \end{titlepage}
    \newpage
}

\begin{document}
    %Titulo del documento. (Portada)
    \portada[Tarea 02]
    %Contenido del documento.
    {\centering\section*{\LARGE El lenguaje C}\vspace{1cm}}
        \subsection*{Historia}
        EL lenguaje C fue creado a mediados de los años 70 a manos de Brian Kernighan y Dennis Ritchie.
        Su primera implementación se realizó sobre el ordenador DEC PDP-11 con sistema operativo
        Unix por Dennis. C es el resultado de un proceso de desarrollo que comenzó con un lenguaje antecesor
        llamado BCPL, el cual influyó en el desarrollo por parte de Khen Thompson de un lenguaje llamado B, el 
        cual es considerado un antecedente de C. Desde su creación, surgieron varias versiones del lenguaje, que 
        incluían unas u otras características, palabras reservadas, etc.
\end{document}